% =====================================================================
%  1bat-ud1-4.tex  —  HORA 4: Formalització d'intervals i inequacions
% =====================================================================
\horatitol{Sessió 4 — Formalització d'intervals i inequacions}%
{La manera «oficial» d'escriure-ho: inequacions, intervals amb parèntesi i claudàtor, i la representació a la recta.}

\subsection{Idea clau}

Ahir van sortir moltes maneres d'escriure «la renda no supera $\num{12600}\eur$»:
«fins a», «com a màxim», $\le\num{12600}$\dots Avui hi posem la notació
\textbf{oficial}: les \textbf{inequacions} i els \textbf{intervals}.

\begin{idea}
\textbf{Inequació:} una expressió amb $<$, $>$, $\le$ o $\ge$ (en lloc de l'$=$).
La seva solució no és un sol valor, sinó \textbf{tot un tram} de la recta.

\textbf{Interval:} la manera d'escriure aquest tram.
\begin{itemize}[leftmargin=1.4em, itemsep=2pt]
  \item \textbf{Claudàtor} $[\,]$: l'extrem \textbf{hi entra} (inclòs, $\le$ o $\ge$).
        A la recta es marca amb un punt \textbf{ple} $\bullet$.
  \item \textbf{Parèntesi} $(\,)$: l'extrem \textbf{no hi entra} (exclòs, $<$ o $>$).
        A la recta es marca amb un punt \textbf{buit} $\circ$.
  \item L'infinit $\infty$ sempre va amb \textbf{parèntesi}: mai s'arriba a l'infinit.
\end{itemize}
\textbf{Exemples de traducció:}
\[
x\le 5 \;\Leftrightarrow\; (-\infty,\,5], \qquad
x>2 \;\Leftrightarrow\; (2,\,+\infty), \qquad
2\le x<7 \;\Leftrightarrow\; [2,\,7).
\]
\end{idea}

\subsection{Exemples}

\begin{exemple}[de la inequació a l'interval i a la recta]
La condició del bo social ordinari era $\text{renda}\le\num{12600}$. Si en diem $x$ a
la renda, $x\le\num{12600}$. En interval és $(-\infty,\,\num{12600}]$ (claudàtor,
perquè el límit hi entra). A la recta:

\begin{center}
\begin{tikzpicture}[>={Stealth[length=2mm]}]
  \draw[->, udgrey] (-4.5,0) -- (4.5,0);
  \draw[udblue, line width=1.4pt] (-4.5,0) -- (2,0);
  \fill[udblue] (2,0) circle (2.4pt);
  \node[below, font=\scriptsize, udgrey] at (2,-0.12) {$\num{12600}$};
  \node[above, udblue, font=\scriptsize] at (-1.5,0.12) {$x\le\num{12600}$};
\end{tikzpicture}
\end{center}

\noindent El punt és \textbf{ple} perquè una renda \emph{exactament} igual a
$\num{12600}\eur$ encara dóna dret.
\end{exemple}

\begin{exemple}[un extrem obert]
«Per accedir cal tenir \textbf{més de} $18$ anys»: si $x$ és l'edat, $x>18$. En
interval, $(18,\,+\infty)$ (parèntesi, perquè el $18$ \emph{no} hi entra). A la recta:

\begin{center}
\begin{tikzpicture}[>={Stealth[length=2mm]}]
  \draw[->, udgrey] (-4.5,0) -- (4.5,0);
  \draw[udblue, line width=1.4pt] (0,0) -- (4.5,0);
  \draw[udblue, fill=white, line width=1.2pt] (0,0) circle (2.4pt);
  \node[below, font=\scriptsize, udgrey] at (0,-0.14) {$18$};
  \node[above, udblue, font=\scriptsize] at (2.2,0.12) {$x>18$};
\end{tikzpicture}
\end{center}

\noindent El punt és \textbf{buit}: tenir exactament $18$ anys no compleix «més de $18$».
\end{exemple}

\subsection{Exercicis resolts}

\begin{exresolt}[escriure interval i recta]
Escriu en forma d'interval i representa a la recta: (a) $x\ge -1$; \ (b) $x<3$.
\solucio
\textbf{(a)} $x\ge -1$ inclou el $-1$: interval $[-1,\,+\infty)$, punt ple a $-1$.

\textbf{(b)} $x<3$ exclou el $3$: interval $(-\infty,\,3)$, punt buit a $3$.

\begin{center}
\begin{tikzpicture}[>={Stealth[length=2mm]}]
  \draw[->, udgrey] (-4,0) -- (4,0);
  \draw[udblue, line width=1.4pt] (-1,0) -- (4,0);
  \fill[udblue] (-1,0) circle (2.2pt);
  \node[below, font=\scriptsize, udgrey] at (-1,-0.12) {$-1$};
  \node[above, udblue, font=\scriptsize] at (1.4,0.12) {(a) $x\ge -1$};
\end{tikzpicture}\\[6pt]
\begin{tikzpicture}[>={Stealth[length=2mm]}]
  \draw[->, udgrey] (-4,0) -- (4,0);
  \draw[udblue, line width=1.4pt] (-4,0) -- (3,0);
  \draw[udblue, fill=white, line width=1.2pt] (3,0) circle (2.2pt);
  \node[below, font=\scriptsize, udgrey] at (3,-0.12) {$3$};
  \node[above, udblue, font=\scriptsize] at (-0.5,0.12) {(b) $x<3$};
\end{tikzpicture}
\end{center}
\resposta{(a) $[-1,+\infty)$, punt ple; (b) $(-\infty,3)$, punt buit.}
\end{exresolt}

\begin{exresolt}[dos límits alhora]
Un programa juvenil admet persones de $16$ anys o més, però menors de $18$. Escriu la
condició com a inequació doble i com a interval.
\solucio
«$16$ anys o més» és $x\ge 16$ (inclòs); «menors de $18$» és $x<18$ (exclòs). Totes
dues alhora:
\[
16\le x<18.
\]
En interval: $[16,\,18)$ (claudàtor a $16$, parèntesi a $18$).
\resposta{$16\le x<18$, és a dir l'interval $[16,18)$.}
\end{exresolt}

\subsection{Exercicis per fer}

\begin{exercicis}
\begin{exlist}
  \item Escriu en forma d'interval:
        \begin{enumerate}[label=\alph*), leftmargin=1.6em, topsep=2pt]
          \item $x\le 4$ \hfill \item $x>-2$ \hfill \item $x\ge 0$ \hfill \item $-3\le x\le 5$
        \end{enumerate}
  \item Escriu com a inequació (o inequació doble) cada interval:
        \begin{enumerate}[label=\alph*), leftmargin=1.6em, topsep=2pt]
          \item $(-\infty,\,7)$ \hfill \item $[2,\,+\infty)$ \hfill \item $(-1,\,4]$
        \end{enumerate}
  \item Representa a la recta real (punt ple o buit segons l'extrem):
        \begin{enumerate}[label=\alph*), leftmargin=1.6em, topsep=2pt]
          \item $x\ge 2$ \hfill \item $x<-1$ \hfill \item $0\le x<3$
        \end{enumerate}
  \item Torna al cas del bo social. Escriu, ara amb \textbf{notació formal} (inequació
        \emph{i} interval), la condició de:
        \begin{enumerate}[label=\alph*), leftmargin=1.6em, topsep=2pt]
          \item família ordinària (llindar $\num{12600}\eur$)
          \item família nombrosa (llindar $\num{16800}\eur$)
        \end{enumerate}
  \item \textbf{(Parèntesi o claudàtor?)} Per a cada cas, digues si l'extrem va amb
        parèntesi o claudàtor i per què:
        \begin{enumerate}[label=\alph*), leftmargin=1.6em, topsep=2pt]
          \item «com a màxim $30$ persones»
          \item «més de $1000\eur$»
        \end{enumerate}
  \item \textbf{(Raonament)} Per què l'infinit s'escriu sempre amb parèntesi i mai amb
        claudàtor? Explica-ho amb les teves paraules.
\end{exlist}
\end{exercicis}

% ---------------------------------------------------------------------
\fullsolucions{Sessió 4}
\begin{solucions}
\begin{sollist}
  \item (a) $(-\infty,4]$; \quad (b) $(-2,+\infty)$; \quad (c) $[0,+\infty)$; \quad
        (d) $[-3,5]$.
  \item (a) $x<7$; \quad (b) $x\ge 2$; \quad (c) $-1<x\le 4$.
  \item (a) punt ple a $2$, ratlla cap a la dreta; \quad (b) punt buit a $-1$, ratlla
        cap a l'esquerra; \quad (c) punt ple a $0$, punt buit a $3$, ratlla entremig.
  \item (a) $x\le\num{12600}$, interval $(-\infty,\,\num{12600}]$; \quad
        (b) $x\le\num{16800}$, interval $(-\infty,\,\num{16800}]$.
  \item (a) claudàtor: «com a màxim $30$» inclou el $30$ ($\le$); \quad
        (b) parèntesi: «més de $1000$» exclou el $1000$ ($>$).
  \item Perquè l'infinit no és un nombre concret al qual s'arribi: és una direcció
        «sense final». Com que mai no s'assoleix, l'extrem no es pot incloure, i per
        això sempre va amb parèntesi.
\end{sollist}
\end{solucions}
